\begin{tabular}{p{0.30\linewidth}p{0.42\linewidth}p{0.20\linewidth}}
\toprule
主張 & 測った量 & 結果\\
\midrule
(1) 線形方程式 $=$ ランダムウォーク & 解とモンテカルロ（2万回）の差の最大（標準誤差の倍数） & 0.9\\
(2) エネルギー最小・キルヒホッフ則 & 自由節点の電流の収支の最大残差；摂動200回がすべて増加 & 2e-14；はい\\
(3) 尺度の共変性 $h(cW;\lambda,c\gamma)=h(W;\lambda,\gamma)$ & 最大差 & 2e-15\\
(4) 三つの行き先の確率の和 $=1$ & 1からのずれの最大 & 2e-15\\
(5) クラスごとの到達確率の和 $=h$ & 最大差 & 2e-15\\
(6) 出口分解（等式），$h\le\pi_Q\le\varphi$ & 分解の誤差の最大；10の塊で$h\le\pi_Q$，$\pi_Q\le\varphi$ & 6e-17；すべて成立\\
(7) 捕まる確率の下界 & 8クラスで$h\ge\kappa\ge\varphi$，$\kappa\ge\rho\omega$ & すべて成立\\
(8) 分離（系1）& 密な未知の塊（種は各塊の5枚に1枚）の種でない160枚：違反数，上界だけで$p\le0.05$が保証される割合（実際）；全手続きの種でない未知154枚での違反数 & 0，94\%（100\%）；0\\
(9) 1次の極限 $(1-h)/\varepsilon\to\mu\mathbb{E}N+\nu\mathbb{E}T$ & 相対誤差の最大（$\varepsilon{=}10^{-2}$，$10^{-4}$）& 0.77，0.031\\
(10) LGCの$u^{+}\to\sqrt{\delta}$（$\lambda\to1$）& Spearman$(u^{+},\sqrt\delta)$（$\lambda{=}0.9\to0.9999$）& 0.55$\to$0.9998\\
(11) 有効性（定理2），400回 & $\alpha{=}0.05/0.1/0.2$でのIDの誤警報率；BH($q{=}0.2$)の偽発見割合 & 0.050/0.101/0.201；0.155\\
\quad 対照：種を$\mathcal{C}$で認定（置換不変でない）& 同上 & 0.054/0.104/0.202；0.175\\
\bottomrule
\end{tabular}
